\documentclass[11pt]{amsart}
\usepackage[T1]{fontenc}
\usepackage[utf8]{inputenc}
\usepackage{amssymb}
\usepackage[margin=1.05in]{geometry}
\usepackage[hidelinks]{hyperref}
\newcommand{\stamptime}{\the\numexpr\time/60\relax:\ifnum\numexpr\time-60*(\time/60)\relax<10 0\fi\the\numexpr\time-60*(\time/60)\relax}
\AddToHook{shipout/foreground}{\put(0.5\paperwidth,-\paperheight+0.55cm){\makebox[0pt]{\tiny\ttfamily unit/doc/\jobname.tex \ --- \ compiled \today, \stamptime}}}
\newtheorem{theorem}{Theorem}
\newtheorem{proposition}[theorem]{Proposition}
\newtheorem{lemma}[theorem]{Lemma}
\newtheorem{corollary}[theorem]{Corollary}
\theoremstyle{definition}
\newtheorem{definition}[theorem]{Definition}
\newtheorem{obligation}[theorem]{Obligation}
\theoremstyle{remark}
\newtheorem{remark}[theorem]{Remark}
\newcommand{\Om}{\Omega}
\newcommand{\Map}{\mathrm{Map}_*}
\newcommand{\Hom}{\mathrm{Hom}}
\newcommand{\Fib}{\mathrm{Fib}}
\newcommand{\Z}{\mathbb{Z}}
\newcommand{\tr}[2]{\|#1\|_{#2}}
\newcommand{\cls}[1]{[#1]}
\newcommand{\hc}[2]{[#1,#2]_*}
\newcommand{\code}[1]{\texttt{#1}}
\title{$\pi_3(S^2)=\Z$ by a Postnikov universal property:\\ informal proof, status of the Unit formalisation}
\date{28 October 2026}
\begin{document}
\maketitle

\noindent\textbf{Status.}
This note answers \code{doc/pi3-s2-postnikov-handoff.txt}. The checked
code is \code{sag/pi3-s2-postnikov} (\code{./segc sag/pi3-s2-postnikov
all}: 0 failed; a fresh process from the directory: 2877
declarations, 0 failed, 1.1\,s). \emph{Update (28 Oct 2026): $\pi_3(S^2)\cong\Z$
is now proved from the sphere axiom alone,
\code{Presentation.Final.pi3-S2} (71 modules, 3135 declarations, 0
failed), through the mapping theorem for targets $\Fib k$
(\code{doc/mapping-fib.pdf}) and a Postnikov presentation of
$\tr{S^2}3$ (\code{doc/postnikov-presentation.pdf}); the text below
records the state before.} \textbf{The theorem $\pi_3(S^2)=\Z$ is
not proved unconditionally.} What is checked is:
\begin{itemize}
\item[(A)] the sphere axiom, its pointed universal property
  $\Map(S^2,X)\simeq\Om^2X$ with naturality; pointed homotopy classes
  $\hc XY$ as \emph{sets}, with composition and the category laws;
  $\pi_3$ as a group, a functor on classes;
\item[(B)] the universal-property argument, \emph{conditional} on two
  explicit interfaces (a 3-truncation of $S^2$; the representing property
  of $P$), concluding $\pi_3(S^2)\cong\pi_3(P)$ as groups;
\item[(C, objects)] for every Postnikov triple $(A,B,k)$, the realization
  $\Fib k$ is a simply connected 3-type and $\pi_3(\Fib k)\cong B$ as
  groups, unconditionally;
\item[(A$'$)] finite truncations $\tr X{m}$ from univalence and
  replacement, and an inhabitant of the interface \code{Trunc3} for
  $S^2$, with the proof that $S^2$ is simply connected
  (\code{doc/finite-truncation.pdf}; \code{SphereTruncation.SphereTrunc.trunc3});
\item[(H)] the cohomology input, unconditionally: the cup square
  $\kappa$ (\code{CupSquare.kappa}, \code{doc/cup-square.pdf}),
  (H3) (\code{H3Proof.h3}, \code{doc/circle-torsors.pdf}) and (H4) for
  this $\kappa$, every abelian $B$ (\code{H4Proof.h4},
  \code{doc/h4.pdf});
\item[(E, conditional)] for any $\kappa:K(\Z,2)\to_*K(\Z,4)$ and the
  mapping theorem for $P=\Fib\kappa$, $\pi_3(S^2)\cong\Z$ as groups
  (\code{Pi3Sphere.Conditional.pi3-S2-T}; the 3-truncation is the
  constructed one).
\end{itemize}
Open: the mapping theorem for the source $P$. The cohomology input
(the construction of $\kappa$ with (H4), and (H3); stated as types in
\code{KZ2Cohomology}, not assumed) was open when this note was first
written. \emph{Update (28 Oct 2026):} (H3) is proved,
\code{H3Proof.h3 : KZ2Cohomology.H3}, from the circle property of $K(\Z,1)$
(\code{CircleProperty}; informal proof \code{doc/circle-torsors.pdf});
$\kappa$ is constructed as the cup square (\code{CupSquare.kappa},
\code{doc/cup-square.pdf}) and (H4) is proved for it,
\code{H4Proof.h4 : KZ2Cohomology.H4 CupSquare.kappa} (\code{doc/h4.pdf}). Sections~\ref{sec:open} and~\ref{sec:audit} give
the precise obligations. No Agda or other external source was inspected
during this work; the references of the handoff were not consulted, and
nothing here depends on them.

The formal files were written in the order of the milestones, before this
note; the note is the informal proof that the files now follow,
declaration by declaration (\S\ref{sec:corr}), in the sense of
\code{seg/WRITING.md}.

\section{Conventions}

Unit's indexing (\code{HLevels}, \code{Connected}): \code{isTrunc n X} is
``$X$ is $(n-2)$-truncated'', \code{isConn n X} is ``$X$ is
$(n-2)$-connected''. So a \emph{3-type} is \code{isTrunc 5} and
\emph{simply connected} is \code{isConn 3}. For Eilenberg--MacLane spaces
we use \code{EMMaps.KK[A].K m}, which is $K(A,m)$ with $K(A,0)=A$; the
offset of \code{EM.EMSpace.K} ($K\,n=K(A,n+1)$) does not appear in the new
files. $\Om X$ is $(x_0=x_0,\mathsf{refl})$; $\Om^2X$ is \code{O2 X x0}.
$\pi_3(X)$ is $\pi_1(\Om^2X)$ (\code{PiOne}: loops modulo mere equality,
composition of the outer loop). $\hc XY$ is the set truncation of
$\Map(X,Y)$, i.e.\ its quotient by mere equality of \emph{pointed} maps
(\code{Quotient}, from replacement).

\section{The sphere (milestone A)}

\begin{definition}[the authorised axiom, \code{SphereAxiom.Sphere}]
A type $S^2$ with $\mathsf{base}$ and $\mathsf{surf}:\Om^2(S^2,\mathsf{base})$
such that for every $X$ in the same universe the canonical map
\[\mathrm{ev}_2:(S^2\to X)\to\textstyle\sum_{x:X}\Om^2(X,x),\qquad
  f\mapsto(f\,\mathsf{base},\ \mathsf{ap}(\mathsf{ap}f)\,\mathsf{surf})\]
is an equivalence. ($\mathsf{ap}f\,\mathsf{refl}$ is $\mathsf{refl}$ by
the computation rule of $J$, so the second component is a 2-loop at
$f\,\mathsf{base}$.)
\end{definition}

This is recursion with its uniqueness (the non-dependent universal
property of the HIT), stated through the one canonical map; it mentions
no homotopy group. The axiom is the parameter telescope of the segment
\code{Sphere}, not a parameter of \code{Theory}.

\begin{lemma}[\code{evPt}, \code{pointed-up}, \code{evPt-nat}]\label{lem:sphere}
Let $\mathrm{ev}_*(F)=(\Om^2F)(\mathsf{surf})$ for $F:S^2\to_*X$. Then
$\mathrm{ev}_*:\Map(S^2,X)\to\Om^2X$ is an equivalence, and
$\mathrm{ev}_*(G\circ F)=(\Om^2G)(\mathrm{ev}_*F)$.
\end{lemma}
\begin{proof}
\emph{Step 1} (\code{Refl.eq}). For $F=(f,\mathsf{refl})$ the looping $J=\Om F$
and $K=(\mathsf{ap}f,\mathsf{refl})$ are homotopic by
$H\,q:J\,q=\mathsf{refl}\cdot\mathsf{ap}f\,q=\mathsf{ap}f\,q$
(\code{loop-ap}, \code{trans-refl-l}). By \code{LoopConj.loop-homotopy},
$(\Om K)\,\mathsf{surf}=c^{-1}\,((\Om J)\,\mathsf{surf})\,c$ with
$c=J_0^{-1}\cdot H(\mathsf{refl})\cdot K_0$, and $c=\mathsf{refl}$ by
computation (\code{c-refl}, stated by \code{refl}). So
$\mathrm{ev}_*F=\mathsf{ap}(\mathsf{ap}f)\,\mathsf{surf}$.
\emph{Step 2} (\code{evPt-subst}). By $J$ on the pointing $p$,
$\mathrm{ev}_*(f,p)=p_*(\mathsf{ap}(\mathsf{ap}f)\,\mathsf{surf})$.
\emph{Step 3} (\code{pointed-up}). $\Map(S^2,X)=\sum_f(f\,\mathsf{base}=x_0)$
is, by $\Sigma$ over the equivalence $\mathrm{ev}_2$ of bases
(\code{sigma-base}: the first component of $\mathrm{ev}_2f$ is
$f\,\mathsf{base}$ by definition), equivalent to
$\sum_{(x,\alpha)}(x=x_0)$, which the transport $((x,\alpha),p)\mapsto
p_*\alpha$ identifies with $\Om^2X$ (\code{Contract}: singleton
contraction). The composite is $\mathrm{ev}_*$ up to Step 2.
\emph{Step 4} (\code{evPt-nat}). $\Om(G\circ F)=\Om G\circ\Om F$ twice
(\code{LoopCalc.loop-comp}).
\end{proof}

\section{Homotopy classes and $\pi_3$ (milestone A)}

\begin{lemma}[\code{HomotopyClasses}]
$\hc XY$ is a set; composition descends in both arguments
(\code{compC}, with \code{comp-beta}: $\cls F;\cls G=\cls{F;G}$); it is
associative and unital (\code{comp-assoc}, \code{comp-idl},
\code{comp-idr}, from the laws for pointed maps). If precomposition with
$F$ is an equivalence of mapping types, precomposition with $\cls F$ is
an equivalence of sets of classes (\code{pre-equiv}, via
\code{classes-equiv}).
\end{lemma}
Equality of classes is \emph{mere} equality of pointed maps
(\code{STr.effective}); nothing is said about the type of homotopies.

\begin{lemma}[\code{HomotopyGroups}]
$\pi_1$ is a functor on pointed maps (\code{ind-comp}, \code{ind-id});
$\Om^2(F;G)=\Om^2F;\Om^2G$ and $\Om^2\mathrm{id}=\mathrm{id}$
(\code{loop2-comp}, \code{loop2-id}); hence $\pi_3$ is a functor on
pointed maps. Homomorphisms form a set (\code{ghom-set}), so
$F\mapsto\pi_3F$ descends to classes (\code{pi3C}), functorially
(\code{pi3C-comp}, \code{pi3C-id}).
\end{lemma}
In particular $\pi_3$ is applied to a class without choosing a
representative; this is what the two-map argument below needs.

\section{The universal-property argument (milestone B, conditional)}

\begin{definition}[\code{SphereUniversal.Interfaces}]
\emph{A 3-truncation of $S^2$} (\code{Trunc3}): a pointed
$\eta:S^2\to_*T$ with $T$ a simply connected 3-type, such that
precomposition with $\eta$ is an equivalence $\Map(T,Y)\simeq\Map(S^2,Y)$
for every 3-type $Y$, and $\pi_3\eta$ is bijective.
\emph{A representing object} (\code{Rep}): a simply connected 3-type $P$
with a class $s\in\hc{S^2}P$ such that, for every simply connected
3-type $X$, $s^*:\hc PX\to\hc{S^2}X$ is a bijection of sets.
\end{definition}

By Lemma~\ref{lem:sphere}, $\hc{S^2}X=\|\Om^2X\|_0=\pi_2X$ (as sets), so
\code{Rep} is ``$\hc PX\cong\pi_2X$ naturally in $X$'', the normalisation
sending $\cls F$ to $\pi_2F$ applied to the element of $\pi_2P$ given by
$s$. The marked object is $(P,s)$; $P$ alone is not an initial object.

\begin{theorem}[\code{Assemble.iso-S2-P}]\label{thm:B}
Given \code{Trunc3} and \code{Rep}, $\pi_3(S^2)\cong\pi_3(P)$ as groups.
\end{theorem}
\begin{proof}
Let $e=\cls\eta$. Precomposition with $e$ is a bijection on classes into
$P$ and into $T$ (\code{etaP}, \code{etaT}: \code{pre-equiv} and
\code{Trunc3.up}); $s^*$ is a bijection on classes into $P$ and $T$
(\code{sP}, \code{sT}: \code{Rep.rep}). Let $u\in\hc TP$ with $e;u=s$
and $t\in\hc PT$ with $s;t=e$ (the inverses at $s$, resp.\ $e$). Then
\[ s;(t;u)=(s;t);u=e;u=s=s;\mathrm{id}_P,\qquad
   e;(u;t)=(e;u);t=s;t=e=e;\mathrm{id}_T,\]
so $t;u=\mathrm{id}_P$ and $u;t=\mathrm{id}_T$ by injectivity
(\code{tu}, \code{ut}). Applying $\pi_3$ on classes, $\pi_3u$ and
$\pi_3t$ are inverse homomorphisms (\code{iso-T-P}). With
$\pi_3\eta$, bijective and hence an isomorphism (\code{gi-of-equiv}:
the inverse of a bijective homomorphism is multiplicative), we get
$\pi_3S^2\cong\pi_3T\cong\pi_3P$.
\end{proof}

Only $X=T$ and $X=P$ are used in \code{Rep.rep}; the general statement is
kept because it is what the mapping theorem gives.

\section{Postnikov triples (milestone C, objects)}

A triple is $(A,B,k)$ with $A,B$ abelian groups and $k:K(A,2)\to_*K(B,4)$;
its realization is the pointed fibre $\Fib k=\sum_{x}(k\,x=\mathsf{pt})$.

\begin{proposition}[\code{PostnikovRealization.Triple}]\label{prop:triple}
$\Fib k$ is a 3-type (\code{trunc5}) and simply connected (\code{conn3}),
and $\pi_3(\Fib k)\cong B$ as groups (\code{pi3-iso}).
\end{proposition}
\begin{proof}
Levels: a $\Sigma$ of a 2-type and of identity types of a 4-type.
Connectedness: $\Fib k$ is 0-connected since every loop of $K(B,4)$ is
merely $\mathsf{refl}$ (\code{Fibres.FibConn}, \code{surj-triv}); its loop
space is $\Fib(\Om k)$ (\code{omega-fib}), 0-connected for the same
reason one level down, so $\Fib k$ is 1-connected
(\code{conn-from-point}).

$\pi_3$: a chain of isomorphisms, each induced by a pointed map whose
looping is an equivalence (\code{LoopIso.LoopEquivIso}: the inverse on
$\pi_1$ descends from the inverse on loops; no choice):
\begin{align*}
\pi_1\Om^2\Fib k
 &\cong \pi_1\Om\Fib(\Om k)        &&\text{($\Om$ of \code{omega-fib}; \code{s1})}\\
 &\cong \pi_1\Fib(\Om^2k)           &&\text{(\code{omega-fib}; \code{s2})}\\
 &\cong \pi_1\Fib(\mathrm{pr}:\Fib(\Om^2k)\to\Om^2K(A,2)) \cong \pi_1\Om^3K(B,4)
   &&\text{(\code{ProjIso}, \code{puppe1}; \code{s3})}\\
 &\cong \pi_1K(B,1)                 &&\text{(loops of \code{KK.e}; \code{s4})}\\
 &\cong B                           &&\text{(the band, a multiplicative bijection; \code{s5}).}
\end{align*}
For \code{s3}: $\Om^2K(A,2)$ is a set, so the projection
$\Fib(\mathrm{pr})\to\Fib(\Om^2k)$ is a subtype inclusion (its fibres are
identity types of a set), whose $\mathsf{ap}$ is an equivalence
(\code{ap-fst-equiv}); and $\Om^3K(B,4)\simeq\Fib(\mathrm{pr})$ is the
Puppe equivalence. This is the boundary $\pi_4K(B,4)\to\pi_3\Fib k$ of the
fibre sequence, read through iterated looping of fibres, preserving the
group operation at each step.
\end{proof}

\begin{corollary}[\code{Pi3Sphere.Conditional.Final.pi3-S2}]
For the sphere axiom, any $\kappa:K(\Z,2)\to_*K(\Z,4)$, a \code{Trunc3}
for $S^2$ and a \code{MapThm} for $P=\Fib\kappa$ (a class
$s\in\hc{S^2}P$ with $s^*$ bijective onto $\hc{S^2}X$ for all simply
connected 3-types $X$): $\pi_3(S^2)\cong\Z$ (\code{Integers.intAb}) as
groups. (Theorem~\ref{thm:B} and Proposition~\ref{prop:triple}.)
\end{corollary}

This is the final signature in conditional form: its parameters
\code{kappa}, \code{TR}, \code{MT} are the open obligations below. They
are explicit parameters of a segment, not Theory assumptions.

\section{Open obligations}\label{sec:open}

\begin{obligation}[finite truncation; feeds \code{Trunc3}. \textbf{DONE}, see \code{doc/finite-truncation.pdf}]
Construct $\tr{S^2}3$ from univalence and replacement, with: the
3-truncation bound, the unit, the universal property for maps into
3-types (pointed form), $\Om^n\tr X{n}\simeq\tr{\Om^nX}0$ for $n=3$ (so
$\pi_3\eta$ is bijective), simple connectivity (which needs that $S^2$ is
1-connected). A route: Rijke's construction of $\tr A{n+1}$ as the image of
$a\mapsto\tr{a=-}n:A\to(A\to U_{\le n})$, by induction from set
truncation (\code{Quotient}); the identity types of the image are
$\tr{a=b}n$ by univalence and the Yoneda lemma for $n$-truncated
families, and the image is small by replacement (\code{Completion}). The
universal property for maps into $(n+1)$-types then has to be proved;
this is the part of the route still to be written out in detail. No
truncation axiom has been added. \emph{Update:} carried out in
\code{ConnectedMaps}, \code{Truncations}, \code{TruncationProps},
\code{SphereTruncation}. The universal property is obtained from the
connectivity of the unit by an elimination lemma for connected maps,
proved with the inductive definition of connectedness, so no join and no
connectivity theorem of the paper is needed.
\end{obligation}

\begin{obligation}[the cohomology input; \code{KZ2Cohomology}. \textbf{DONE}, see \code{doc/cup-square.pdf}, \code{doc/h4.pdf}]
Construct $\kappa$ and prove, for every abelian group $B$,
\[(\mathrm{H4})\quad \Phi_B:\Hom(\Z,B)\to\hc{K(\Z,2)}{K(B,4)},\
  g\mapsto\cls{\kappa;K(g,4)}\ \text{is a bijection},\qquad
  (\mathrm{H3})\quad \hc{K(\Z,2)}{K(B,3)}\ \text{is contractible}.\]
Both are stated as types (\code{H4 kappa}, \code{H3}); neither is assumed.
\emph{Update:} (H3) is now proved (\code{H3Proof.h3}) along the first
reduction below, with the circle property of the $\Z$-torsors proved in
\code{CircleProperty} (\code{doc/circle-torsors.pdf}). (H4) is proved for
the cup square $\kappa$ (\code{H4Proof.h4}, \code{doc/h4.pdf}): the
second reduction below is Proposition~5 there (\code{H4Classes}); the
remaining content is obtained without pushouts and without the
quadratic-map route, as $\Phi_B=\mathrm{cls}\circ\Delta'\circ(g\mapsto
c_2;K(g,4)_*)$, a composite of bijections: every class is the diagonal of
a unique bipointed map (Theorem~A, \code{H4Diagonal}) and $g\mapsto
c_2;K(g,4)_*$ is a bijection onto the bipointed maps (Theorem~B,
\code{H4Bilinear}), for every abelian $B$, $B=\Z/2$ included. The
paragraphs below record the state before that proof. Informal reductions:
\begin{itemize}
\item (H3). Looping is a map $\Map(K(\Z,2),K(B,3))\to\Map(K(\Z,1),K(B,2))$
  whose fibres (the deloopings) are propositions, by DW1 Cor.~10
  ($K(\Z,2)$ is 1-connected, $K(B,3)$ is 3-truncated:
  $2\cdot1+m+2=3$, $m=-1$), so it is an embedding. If $K(\Z,1)$ (the
  $\Z$-torsors) has the universal property of the circle,
  $\Map(K(\Z,1),K(B,2))\simeq\Om K(B,2)\simeq K(B,1)$ is connected. A
  subtype of a connected type containing the base point has one
  component, so $\hc{K(\Z,2)}{K(B,3)}$ is a point. The missing lemma is the
  circle property of $\Z$-torsors (Bezem--Buchholtz--Grayson--Shulman);
  it is not in the library.
\item (H4). The same looping, now with fibres of level $m=0$ (sets), maps
  $\Map(K(\Z,2),K(B,4))$ to $\Map(K(\Z,1),K(B,3))\simeq K(B,2)$, which is
  simply connected; so $\hc{K(\Z,2)}{K(B,4)}$ is the set of deloopings of
  the zero map $K(\Z,1)\to_*K(B,3)$. That this set is $B$, with the
  coefficient changes, is the real content (classically the Hopf
  construction, or $H^4(\mathbb{CP}^\infty)$). I do not have a synthetic
  proof of it from the authorised foundations: the cellular route through
  $\mathbb{CP}^\infty$ needs pushouts, which are not authorised. This is the
  main open mathematical obligation. The quadratic-map route of the
  handoff ($q(-a)=q(a)$, bilinear cross-effect; $\Gamma(\Z)=\Z$,
  generator with $q(1)$ a generator, not $2q(1)$) would equally need a
  comparison theorem with pointed maps $K(A,2)\to K(B,4)$, which is not
  available.
\end{itemize}
\end{obligation}

\begin{obligation}[the mapping theorem for the source $P=\Fib\kappa$; feeds \code{MapThm}]
Informal proof, given (H3), (H4) and the obligation on truncations, for a
target $X=\Fib k$ with $k:K(A,2)\to_*K(B,4)$:
\begin{enumerate}
\item $\Map(P,\Fib k)\simeq\sum_{G:\Map(P,K(A,2))}(G;k=_*\mathrm{const})$
  (universal property of the pointed fibre).
\item Precomposition with $\mathrm{pr}:P\to K(\Z,2)$ is an equivalence
  $\Map(K(\Z,2),K(A,2))\simeq\Map(P,K(A,2))$: the fibres of $\mathrm{pr}$
  are merely $K(\Z,3)$ (2-connected) and $K(A,2)$ is a 2-type (Lemma 4 in
  the pointed form, fibrewise). And $\Map(K(\Z,2),K(A,2))\simeq\Hom(\Z,A)
  \simeq A$ is a set (unique deloopings, DW1 Cor.~11; evaluation at $1$).
\item Hence $\hc P{\Fib k}\cong\sum_{a:A}\pi_0\,\mathrm{Null}(h_a)$ with
  $h_a=\mathrm{pr};K(a,2);k$. By (H4), $K(a,2);k=\kappa;K(g,4)$ merely,
  for a unique $g$, so $h_a$ is merely $\mathrm{pr};\kappa;K(g,4)$, which
  is null by the fibre sequence. A path space in
  $\Map(P,K(B,4))$ that is inhabited has $\pi_0$ a torsor under
  $\hc P{K(B,3)}$ (translation by the additive structure of
  $K(B,4)$, or equivalently path composition); so it is a point once
  $\hc P{K(B,3)}=0$.
\item $\hc P{K(B,3)}=0$: from the low-degree exact sequence of the
  fibration $K(\Z,3)\to P\to K(\Z,2)$,
  $0=H^3(K(\Z,2);B)\to H^3(P;B)\to H^3(K(\Z,3);B)\xrightarrow{\tau}
  H^4(K(\Z,2);B)$, with $\tau=\Phi_B$ injective by (H4). The exactness
  (a synthetic Serre sequence in degrees $\le4$, base 1-connected, fibre
  2-connected) is itself to be proved. \emph{Update (28 Oct 2026):}
  \textbf{proved}, without a Serre sequence, for $P=\Fib\kappa$ and any
  $\kappa$ with (H4): \code{FibreH3.fib-h3}
  (\code{doc/fibre-h3.pdf}, refereed). The restriction of $F$ to the fibre
  deloops to $G$; $F$ gives an explicit null-homotopy of $\kappa;G$ (so
  the connecting map is $\Phi_B$ by construction); $G=K(g,4)$ and (H4)
  give $G=0$; then $F$ factors through $K(\Z,2)$ and (H3) kills it.
\item Naturality in $X$ and the reduction of an arbitrary $X$ in $C$ to a
  presentation $\Fib k$: the statement ``$s^*$ is bijective at $X$'' is a
  proposition, so a mere presentation suffices (no global choice). The
  presentation needs $X\to\tr X2\simeq K(\pi_2X,2)$, the fibre a
  2-connected 3-type $\simeq K(\pi_3X,3)$, and the classification of the
  resulting fibration over a simply connected base by a map to
  $K(\pi_3X,4)$. For Theorem~\ref{thm:B} only $X=P$ (a presentation by
  definition) and $X=T$ are needed.
\end{enumerate}
None of this is formalised. The handoff's triple-morphism formulation
($f,g,H$ with $H$ a pointed homotopy, the $H^3$-ambiguity removed by (H3)
on components, not on homotopies) is equivalent for this source; the
reduction above was preferred because it avoids building the general
morphism theory.
\end{obligation}

\section{Correspondence with the formal declarations}\label{sec:corr}

\begin{tabular}{ll}
Lemma~\ref{lem:sphere} & \code{SphereAxiom}: \code{Sphere.Refl.eq}, \code{evPt-subst}, \code{pointed-up}, \code{evPt-nat}\\
classes & \code{HomotopyClasses}: \code{STr}, \code{HC}, \code{compC}, \code{comp-assoc}, \code{pre-equiv}\\
$\pi_3$ on classes & \code{HomotopyGroups}: \code{pi3}, \code{pi3-map}, \code{pi3C}, \code{pi3C-comp}, \code{GIso}\\
Theorem~\ref{thm:B} & \code{SphereUniversal}: \code{Interfaces}, \code{Assemble.tu}, \code{ut}, \code{iso-S2-P}\\
loop isomorphisms & \code{LoopIso}: \code{loopPE}, \code{LoopEquivIso}, \code{ProjIso}, \code{LoopsSetIso}\\
Proposition~\ref{prop:triple} & \code{PostnikovRealization.Triple}: \code{trunc5}, \code{conn3}, \code{s1}--\code{s5}, \code{pi3-iso}\\
(H3), (H4) as types & \code{KZ2Cohomology}: \code{Phi}, \code{H4}, \code{H3}\\
(H3), (H4) proved & \code{H3Proof.h3}; \code{CupSquare.kappa}, \code{H4Proof.h4}\\
Corollary & \code{Pi3Sphere.Conditional.Final.pi3-S2}
\end{tabular}

\section{Audit}\label{sec:audit}

\emph{Assumptions.} The \code{Theory} parameters of \code{sag/em}
(univalence, images, surjectivity; function extensionality derived).
The sphere axiom as a segment parameter. No other axiom; the checker was
not modified; no choice principle.

\emph{Checking against the text} (\code{seg/WRITING.md}, uniform checker
\code{headlinear}, \code{--opened} and \code{--opens pi3.opens}). Every
declaration checks in under 0.01\,s. \code{c-refl}, \code{conn3},
\code{pi3-iso}, \code{iso-S2-P} open nothing beyond their text. The
extras that remain: (i) \code{Refl.eq}, \code{evPt-nat}: the base point of
a looping $\Om F\,\mathsf{refl}$ computes to $\mathsf{refl}$
(\code{Maps.D}, \code{Kf}, \code{idJR}), a definitional fact used in types;
(ii) \code{tu}, \code{ut}, \code{comp-assoc}: the carrier $\hc XY$ is opened
to the quotient \code{Quot.Q} and its \code{image}; (iii) the $\pi_1$
lemmas open \code{PiOne.grp}/\code{QGrp.G} to \code{Quot.Q}, because
the library states \code{Induced.map-cls} with \code{grp.G} while
\code{Descend} works on \code{Q.Q}. (ii) and (iii) come from the
interfaces of the shared library (\code{Quotient}, \code{QGrp},
\code{PiOne}), which were not edited. Aliases of named notions were
removed (\code{Assemble} uses \code{open TR}, \code{open RP}; \code{Triple}
writes \code{KA.K two}); \code{GIso} is a record with accessors.

\emph{Not done.} $\pi_2$ as a group;
$\Hom(\Z,B)\cong B$; the checks of the handoff on triples
($K(A,2)=(A,0,0)$, $\Fib0\simeq K(A,2)\times K(B,3)$, induced maps of
squares on $\pi_2,\pi_3$, $\Phi_B(0)=0$, $\Phi_\Z(\mathrm{id})=\cls\kappa$).
\end{document}
